Um die Antikommutatorrelation der Dirac-Matrizen zu zeigen, beginnen wir mit der Definition des Antikommutators:

\[
\{\gamma^\mu, \gamma^\nu\} = \gamma^\mu \gamma^\nu + \gamma^\nu \gamma^\mu
\]

Diese Relation ist definiert, um die Clifford-Algebra der Dirac-Matrizen zu erfüllen, die in der Form

\[
\{\gamma^\mu, \gamma^\nu\} = 2 g^{\mu \nu} I
\]

gegeben ist, wobei \( g^{\mu \nu} \) die Metrik des Minkowski-Raums ist und \( I \) die Einheitsmatrix darstellt. Diese Gleichung ist eine fundamentale Eigenschaft der Dirac-Matrizen und wird als gegeben angenommen.

Um nun die zweite Gleichung zu zeigen, verwenden wir die Eigenschaft der Dirac-Matrizen in Bezug auf die Hermitesche Konjugation. Wir wissen, dass die Dirac-Matrizen die Beziehung

\[
\gamma^{0 \dagger} = \gamma^0, \quad \gamma^{i \dagger} = -\gamma^i
\]

erfüllen, wobei \( i \) die räumlichen Indizes darstellt. Um die gegebene Gleichung zu zeigen, betrachten wir die Hermitesche Konjugation von \(\gamma^\mu\):

\[
\gamma^{\mu \dagger} = (\gamma^0 \gamma^\mu \gamma^0)^\dagger
\]

Da die Hermitesche Konjugation die Reihenfolge der Matrizen umkehrt und die Konjugationseigenschaften der Dirac-Matrizen berücksichtigt, erhalten wir:

\[
\gamma^{\mu \dagger} = \gamma^0 (\gamma^\mu)^\dagger \gamma^0
\]

Da \((\gamma^\mu)^\dagger = \gamma^0 \gamma^\mu \gamma^0\) gilt, folgt daraus:

\[
\gamma^{\mu \dagger} = \gamma^0 (\gamma^0 \gamma^\mu \gamma^0) \gamma^0 = \gamma^0 \gamma^\mu \gamma^0
\]

Dies zeigt die geforderte Beziehung.